\subsection{Характеристическое уравнение}

Характеристическое уравнение системы управления в общем виде определяется выражением
\begin{equation}
    \det A(p) = 0,
\end{equation}
где $A(p)$ --- матрица коэффициентов системы уравнений движения модуля поворота.

Корни характеристического уравнения определяют вид решения системы и, в частности, ее устойчивость. Для исследуемой системы матрица $A(p)$ была получена ранее:
\begin{equation}
    A=
    \begin{bmatrix}
        J_{\text{дпр}}p^2+bp+c
        &
        -(bp+c)
        &
        -1
        &
        0
        \\

        -(bp+c)
        &
        J_{\text{м}}p^2+bp+c
        &
        0
        &
        0
        \\

        s_{\text{пр}}p
        &
        0
        &
        \tau p+1
        &
        -r_{\text{пр}}
        \\

        kk_{\text{п}}p
        &
        kk_{\text{и}}
        &
        0
        &
        p
    \end{bmatrix}.
\end{equation}

Исходные данные для расчета:
\[
    J_{\text{д}} = 2 \cdot 10^{-3}, \quad
    J_{\text{м}} = 1, \quad
    b = 4, \quad
    c = 10^3,
\]
\[
    i = 50, \quad
    \tau = 3 \cdot 10^{-3}, \quad
    s = 0.05, \quad
    r = 0.7, \quad
    k = 100.
\]

Приведенные коэффициенты:
\begin{equation}
    J_{\text{дпр}} = J_{\text{д}} i^2 = 5,
\end{equation}

\begin{equation}
    s_{\text{пр}} = s i^2 = 125,
\end{equation}

\begin{equation}
    r_{\text{пр}} = r i = 35.
\end{equation}

Подставляя численные значения параметров в матрицу $A(p)$, получим:
\begin{equation}
    A(p)=
    \begin{bmatrix}
        5p^2+4p+1000
        &
        -(4p+1000)
        &
        -1
        &
        0
        \\

        -(4p+1000)
        &
        p^2+4p+1000
        &
        0
        &
        0
        \\

        125p
        &
        0
        &
        0.003p+1
        &
        -35
        \\

        100k_{\text{п}}p
        &
        100k_{\text{и}}
        &
        0
        &
        p
    \end{bmatrix}.
\end{equation}

Найдем определитель матрицы $A(p)$~--- результат вычислений в Wolfram Mathematica приведён на рисунке~\ref{fig:determinant_wolfram}, код программы~--- на листинге~\ref{lst:determinant}.

\begin{listing}[H]
\caption{Нахождение определителя матрицы $A(p)$}
\vspace{-10pt}
\label{lst:determinant}
\begin{minted}[linenos, numbersep=10pt]{mathematica}
Clear[p, kp, ki];

Jd = 2*10^-3;
Jm = 1;
b = 4;
c = 10^3;
i = 50;
tau = 3*10^-3;
s = 0.05;
r = 0.7;
k = 100;

Jdpr = Jd*i^2;
spr = s*i^2;
rpr = r*i;

A = {{Jdpr*p^2 + b*p + c, -(b*p + c), -1, 0}, {-(b*p + c),
    Jm*p^2 + b*p + c, 0, 0}, {spr*p, 0, tau*p + 1, -rpr}, {k*kp*p,
    k*ki, 0, p}};

detA = Expand[Det[A]];

coeffs = CoefficientList[detA, p];

coeffsDescending = Reverse[coeffs];

detA == 0
\end{minted}
\end{listing}

\screenshot[0.95]{determinant_wolfram}{
Результат вычисления характеристического полинома в системе Wolfram Mathematica
}

\begin{equation}
    \det A(p)
    =
    0.015p^6
    +
    5.072p^5
    +
    167p^4
    +
    (3500k_{\text{п}}+6500)p^3
\end{equation}
\[
    \qquad
    +
    (14000k_{\text{п}}+125000)p^2
    +
    (14000k_{\text{и}}+3500000k_{\text{п}})p
    +
    3500000k_{\text{и}}.
\]

Таким образом, характеристическое уравнение системы имеет вид
\begin{empheq}[box=\fbox]{equation}
\begin{aligned}
    0.015p^6
    &+
    5.072p^5
    +
    167p^4
    +
    (3500k_{\text{п}}+6500)p^3
    \\
    &+
    (14000k_{\text{п}}+125000)p^2
    \\
    &+
    (14000k_{\text{и}}+3500000k_{\text{п}})p
    +
    3500000k_{\text{и}}
    =
    0.
\end{aligned}
\end{empheq}

Обозначим коэффициенты характеристического полинома:
\begin{equation}
    a_0 = 0.015,
\end{equation}

\begin{equation}
    a_1 = 5.072,
\end{equation}

\begin{equation}
    a_2 = 167,
\end{equation}

\begin{equation}
    a_3 = 3500k_{\text{п}}+6500,
\end{equation}

\begin{equation}
    a_4 = 14000k_{\text{п}}+125000,
\end{equation}

\begin{equation}
    a_5 = 14000k_{\text{и}}+3500000k_{\text{п}},
\end{equation}

\begin{equation}
    a_6 = 3500000k_{\text{и}}.
\end{equation}
